\scanpage{76}
% loc: 1–4, top, with polynomial substitutions
\[
\begin{aligned}
 \delta(\underline{xyzt})&=\delta(xy(zt))=Q\bigl(\delta(x),\delta(y),\delta(zt),\delta(xy),\underline{\delta(yzt)},\underline{\delta(ztx)}\bigr)\\
 &\deleted{\ne\delta(x(yz),t)}\\
 &=\delta(x(yz)t)=Q\bigl(\delta(x),\delta(yz),\delta(t),\underline{\delta(xyz)},\underline{\delta(yzt)},\delta(tx)\bigr)\\
 &=\delta((xy)zt)=Q\bigl(\delta(xy),\delta(z),\delta(t),\underline{\delta(xyz)},\delta(zt),\underline{\delta(txy)}\bigr)
\end{aligned}
\]
\margin{\unsure{remplacer par|rempl. par|remplaçons les} Pol. en $\delta(y),\delta(z)\,\unsure{\&|,}\,\delta(t)$}

% loc: 5–7, beneath equations
\diagram{\begin{tikzcd}[ampersand replacement=\&,column sep=small,row sep=small]
 {} \& yzt\arrow[dr,dash] \& \deleted{\unsure{z}}ztx\\
 xyz\arrow[dr,dash] \& {} \& yzt\\
 {} \& xyz \& txy
\end{tikzcd}}

% loc: 8–11
On trouve trois expressions polynômiales différentes d'une $=$ quantité, d'où 2 relations (\unsure{prob\textsuperscript{t}|probl\textsuperscript{t}|quasi !} \unsure{équivalentes|arguables}) entre les dix quantités
\[
 \delta(x),\delta(y),\delta(z),\delta(t),\quad\delta(xy)\ \delta(xz)\ \delta(xt)\ \delta(yz)\ \delta(yt)\ \delta(zt)
\]

% loc: 12–14
Exprimer que si $\tau$ normalise $H\subset G$, $H$ \uline{comm.}, alors $h\in H$,
\[
 \boxed{\delta(\tau h)\overset{!}{=}\delta(\tau)}
\]
i.e. les $\tau^n(h)\inserted{\overset{\mathrm{déf}}{=}\tau^nh\tau^{-n}}$ commutent à $h$ ($n\in\mathbb Z$), alors on a la formule encadrée. On aimerait l'écrire comme une ``identité'' inconditionnelle

% loc: left, oblique beside 12–17
% slanted
\margin{C'est \uline{idiot}, \unsure{comme on voit|comme on sait} en prenant $\tau\in H$ -- il faut supposer de plus $\tau$ \unsure{ou} ``$\tau$ moins \unsure{un}'' opère de façon \unsure{non triviale|trivialement} sur $H$, i.e. qque chose du goût $\ldots$}

% loc: 15–17
\[
 \delta(\tau h)-\delta(\tau)\overset{?}{=}
 \text{expression polynomiale}\ \inserted{(\text{à coeff}\in\mathbb Z)\ \deleted{\text{\unsure{système constant|système constante}}}}\ \text{en les}
\]
\[
 \delta\bigl(\underbrace{\{\tau^n(h),h\}}_{\text{commutateur}}\bigr)
 \ \text{et }\delta(\tau),\delta(\tau h),\delta(h)
\]
\inserted{dont tous les termes contiennent un $\delta(\{\cdot,\cdot\})$}

% loc: 18–19
(peut-être suffit-il de
\[
 \underbrace{\{\tau h\tau^{-1},h\}}_{\tau h\tau^{-1}h\tau h^{-1}\tau^{-1}h^{-1}}
 \ \text{et}\quad
 \underbrace{\{\tau^{-1}h\tau,h\}}_{\tau^{-1}h\tau h\tau^{-1}h^{-1}\tau h^{-1}}
\]

% loc: 20, bottom edge
A-t-on au moins
\[
 \boxed{\delta(\tau.\underbrace{\tau.(h)h^{-1}}_{\{\tau,h\}})\overset{?}{=}\delta(\tau)}
 \quad(\text{si }\tau\ \text{normalise }H\ \text{comm, }h\in H)
\]
% archive sheet number: 75

